\documentclass[11pt,oneside]{article}

\usepackage[T1]{fontenc}
\usepackage{lmodern}
\usepackage{microtype}
\usepackage{geometry}
\geometry{margin=1in,headheight=14pt}
\usepackage{amsmath,amssymb,amsthm,mathtools}
\usepackage{booktabs,longtable,tabularx,array,multirow}
\usepackage{enumitem}
\usepackage{xcolor}
\usepackage{hyperref}
\usepackage[nameinlink,noabbrev]{cleveref}
\usepackage{listings}
\usepackage{fancyhdr}
\usepackage{titlesec}
\usepackage{needspace}
\usepackage{url}

\definecolor{leanblue}{RGB}{24,91,145}
\definecolor{newgreen}{RGB}{15,105,68}
\definecolor{computeviolet}{RGB}{98,55,145}
\definecolor{rulegray}{RGB}{92,101,111}
\definecolor{codegray}{RGB}{248,248,248}

\hypersetup{
  colorlinks=true,
  linkcolor=leanblue,
  citecolor=newgreen,
  urlcolor=leanblue,
  pdftitle={The Alaoglu--Erdos Integer-Exponent Conjecture: Combined Report},
  pdfauthor={The ProveIt project},
  pdfsubject={Machine-checked partial results, exact conditional structure, and research directions}
}
\urlstyle{same}

\pagestyle{fancy}
\fancyhf{}
\lhead{Alaoglu--Erd\H{o}s integer-exponent conjecture}
\rhead{Combined report}
\cfoot{\thepage}
\renewcommand{\headrulewidth}{0.4pt}

\titleformat{\section}{\Large\bfseries}{\thesection}{0.75em}{}
\titleformat{\subsection}{\large\bfseries}{\thesubsection}{0.75em}{}
\titleformat{\subsubsection}{\normalsize\bfseries}{\thesubsubsection}{0.75em}{}

\setlist[itemize]{leftmargin=2em,itemsep=0.3em,topsep=0.4em}
\setlist[enumerate]{leftmargin=2.2em,itemsep=0.3em,topsep=0.4em}
\setlength{\parindent}{0pt}
\setlength{\parskip}{0.55em}
\setlength{\emergencystretch}{3.5em}

\newtheorem{theorem}{Theorem}[section]
\newtheorem{proposition}[theorem]{Proposition}
\newtheorem{lemma}[theorem]{Lemma}
\newtheorem{corollary}[theorem]{Corollary}
\newtheorem{conjecture}[theorem]{Conjecture}
\newtheorem{question}[theorem]{Question}
\theoremstyle{definition}
\newtheorem{definition}[theorem]{Definition}
\newtheorem{observation}[theorem]{Observation}
\theoremstyle{remark}
\newtheorem{remark}[theorem]{Remark}

\newcommand{\N}{\mathbb N}
\newcommand{\Nzero}{\mathbb N_0}
\newcommand{\Z}{\mathbb Z}
\newcommand{\Q}{\mathbb Q}
\newcommand{\R}{\mathbb R}
\newcommand{\C}{\mathbb C}
\newcommand{\fract}[1]{\left\{#1\right\}}
\newcommand{\distZ}[1]{\left\lVert #1\right\rVert_{\Z}}
\newcommand{\vpp}[2]{v_{#1}\!\left(#2\right)}
\newcommand{\thetaq}{\vartheta}
\newcommand{\Sol}{\mathcal S}
\newcommand{\Cand}{\mathcal C}
\newcommand{\Pair}{\mathcal P}
\newcommand{\Zthree}{\Z[1/3]}
\newcommand{\clog}{\operatorname{clog}}
\newcommand{\sha}[1]{\texttt{\detokenize{#1}}}
\newcommand{\lean}[1]{\texttt{\detokenize{#1}}}
\newcommand{\commit}[1]{\texttt{#1}}
\newcommand{\statuslean}{\textcolor{leanblue}{\textbf{[kernel-verified Lean]}}}
\newcommand{\statusnew}{\textcolor{newgreen}{\textbf{[paper proof in this report]}}}
\newcommand{\statuscomp}{\textcolor{computeviolet}{\textbf{[interval computation]}}}
\newcommand{\statusknown}{\textcolor{rulegray}{\textbf{[classical theorem]}}}

\newenvironment{statusbox}[1]
{\begin{center}\begin{minipage}{0.94\textwidth}\raggedright\hbadness=10000\hrule\vspace{0.45em}\textbf{Status:} #1\par\vspace{0.35em}}
{\vspace{0.35em}\hrule\end{minipage}\end{center}}

\lstdefinestyle{pythonstyle}{
  language=Python,
  basicstyle=\ttfamily\footnotesize,
  backgroundcolor=\color{codegray},
  frame=single,
  rulecolor=\color{rulegray},
  numbers=left,
  numberstyle=\tiny\color{rulegray},
  stepnumber=1,
  numbersep=8pt,
  showstringspaces=false,
  breaklines=true,
  columns=fullflexible,
  keepspaces=true,
  tabsize=4,
  captionpos=b
}
\lstdefinestyle{leanstyle}{
  basicstyle=\ttfamily\footnotesize,
  backgroundcolor=\color{codegray},
  frame=single,
  rulecolor=\color{rulegray},
  showstringspaces=false,
  breaklines=true,
  columns=fullflexible,
  keepspaces=true,
  tabsize=2,
  captionpos=b
}

\begin{document}
\pagenumbering{gobble}

\begin{titlepage}
\centering
\vspace*{1.0cm}
{\Huge\bfseries The Alaoglu--Erd\H{o}s\par}
\vspace{0.25cm}
{\Huge\bfseries Integer-Exponent Conjecture\par}
\vspace{0.65cm}
{\LARGE Machine-Checked Partial Results,\par}
\vspace{0.15cm}
{\LARGE Exact Conditional Structure,\par}
\vspace{0.15cm}
{\LARGE and Research Directions\par}
\vspace{1.05cm}
{\large A combined survey and research report on the corpus in\par}
\vspace{0.15cm}
{\large \url{https://github.com/VladimirReshetnikov/ProveIt}\par}
\vfill
\begin{minipage}{0.88\textwidth}
\small
\textbf{Bottom line.}
No unconditional proof of the two-base conjecture is obtained, and the reduction to the four
exponentials conjecture explains why none should be expected from routine combination of the
present methods.  The combined corpus does establish, in kernel-verified Lean: a zero-free
region $2^x<4096$ (any nonintegral solution has $x\ge12$); transcendence of every hypothetical
nonintegral solution; a common odd core $w$ with every candidate of the form $2^i w^j$;
logarithmic counting and a linear dyadic-block bound; local progression obstructions with
density zero; explicit Diophantine lower bounds; the classification of further integral bases
as exactly the $3$-smooth numbers; and unconditional equivalences between the conjecture and
fractional-part gap statements.  At the paper level this report sharpens the conditional
picture to an exact rigidity theorem: if one counterexample exists, there is a unique least
nonintegral solution $\beta$ and
\[
   \Sol=\{n+k\beta:n,k\in\Nzero\},
\]
with exact counts in every unit interval and Weyl equidistribution of fractional parts.  An
independent interval computation excludes all nonintegral candidates with $2^x<2^{20}$; that
computation is explicitly \emph{not} a Lean kernel proof.  The irreducible remainder is
isolated as a radical statement: no odd $u>1$ has $u^{\log_2 3}\in\Z[1/3]$.
\end{minipage}
\vfill
{\normalsize This document merges, updates, and extends two earlier AI-assisted reports\par
(Claude Fable~5, Anthropic; GPT-5.6 Pro, OpenAI), against the repository state of\par
commit \commit{716617d72}.  All Lean identifiers cited below are the authoritative
statements.\par}
\vspace{0.6cm}
{\large August 20, 2026\par}
\end{titlepage}
\pagenumbering{roman}

\begin{abstract}
The Alaoglu--Erd\H{o}s conjecture asks whether a real $x$ satisfying $2^x,3^x\in\Z$ must be
an integer.  It is a particularly concrete unresolved instance of the four exponentials
problem: the three-base analogue with $5^x\in\Z$ added is a theorem, formalized in the
repository by an interpolation-determinant argument, while the two-base statement sits exactly
at the open $2\times2$ transcendence threshold.

This report is the merged and expanded successor of two documents: a survey of the
machine-checked partial results and an independent research report containing an audit and new
conditional structure theory.  It serves four purposes.  First, it catalogues the
kernel-verified Lean corpus --- zero-free regions through $2^x<4096$, Gelfond--Schneider
transcendence restrictions, multiplicative rank at most two with a common odd core,
progression obstructions and zero density, shrinking-target Diophantine bounds, the $3$-smooth
classification of auxiliary bases, the four-exponentials reduction, and unconditional
equivalences between the conjecture and fractional-part gap statements --- with the precise
Lean identifiers and their verification status.  Second, it develops the exact conditional
structure: if a counterexample exists, a canonical primitive odd core $w$, rational-power
index $d$, and least nonintegral solution $\beta=a+d\log_2w$ exist, and the full solution set
is the free additive monoid $\Nzero+\Nzero\beta$; exact dyadic-block counts and blockwise Weyl
equidistribution follow.  Third, it records a sharper family of finite-difference obstructions
for progressions of any length, and an independent 40-digit interval verification excluding
all nonintegral candidates below $2^{20}$ under a stated software trust boundary.  Fourth, it
consolidates the research program: the remaining obstruction is equivalent to the radical
statement that $u^{\log_2 3}\notin\Z[1/3]$ for every odd $u>1$, and the report prioritizes
formalization targets and genuinely new attack surfaces around that statement.
\end{abstract}

\tableofcontents
\clearpage
\pagenumbering{arabic}

\section{Executive summary}
\label{sec:executive}

\subsection{What is proved, and at what level of trust}

The full conjecture remains open.  Several routes were pushed to their exact failure points:
the three-base determinant method does not shrink to two bases without proving a case of the
four exponentials conjecture; equidistribution of multiples does not defeat exponentially
shrinking integrality gaps; local curvature obstructions cannot reach the sparse global
structure; and every successful reduction terminates at the same transcendence-theoretic
wall, now isolated in Conjecture~\ref{conj:radical-target}.

To keep unlike kinds of evidence separate, every important statement in this report carries
one of four labels.

\begin{description}[leftmargin=3.9cm,style=nextline]
\item[\statuslean]
A proof source is committed in the repository, lies on the audited
\lean{ExponentialIdentities.lean} import surface, and has been accepted by the Lean kernel in
a full build of the cited modules (toolchain 4.31/4.32 series, Mathlib pinned by the
repository manifest).  No \lean{sorry}, no extra axioms, no \lean{native_decide}.
\item[\statusnew]
A conventional mathematical proof is supplied in this document.  It is intended to be
suitable for formalization, but no claim of kernel checking is made.
\item[\statuscomp]
A finite statement was checked by outward-rounded interval arithmetic.  The script,
parameters, worst margins, and manifest hashes are recorded in
Section~\ref{sec:finite-computation} and Appendix~\ref{app:interval-output}.  The trusted base includes CPython and
\lean{mpmath}; this is not interchangeable with a Lean certificate.
\item[\statusknown]
A classical theorem quoted from the literature.
\end{description}

An earlier audit, performed against commit \commit{423ac24}, correctly observed that three
modules cited by the survey (\lean{FiniteCheck4096}, \lean{OddCore},
\lean{FractionalPartDensity}) were absent from that commit tree.  That discrepancy is
resolved: the three modules were subsequently committed in \commit{98646b455} and are
kernel-verified; the survey's temporary pending-verification remark was removed in
\commit{716617d72}.  Statements from these modules are therefore labeled \statuslean{} below.

\subsection{Headline results}

\begin{enumerate}
\item \textbf{Verified zero-free region.}  If $2^x,3^x\in\Z$ and $2^x<4096$, then $x\in\Z$;
equivalently every nonintegral solution has $x\ge12$.  \statuslean{}  An independent interval
scan extends the finite barrier to $2^x<2^{20}$, i.e.\ $x>20$.  \statuscomp
\item \textbf{Transcendence.}  Every nonintegral solution is transcendental; for $x$ with
$2^x\in\Z$, algebraicity of $x$ is \emph{equivalent} to integrality.  $\thetaq=\log_23$ is
transcendental.  \statuslean
\item \textbf{Structure.}  All candidate values $2^x$ lie in $\{2^iw^j\}$ for one odd $w>1$;
candidate counts below $N$ are at most $\clog_2N\cdot\clog_3N$, and each dyadic block
$[2^T,2^{T+1})$ holds at most $\clog_3(2^{T+1})$ candidates.  \statuslean{}  Conditional on a
counterexample, this sharpens to an exact free-monoid classification
$\Sol=\Nzero+\Nzero\beta$ with a canonical least generator, exact block counts
$\lfloor(N+1)/\beta\rfloor$, and blockwise Weyl equidistribution.  The generator theorem
itself is kernel-verified (\lean{exists_primitive_generator} in
\lean{PrimitiveGenerator}); the exact counts and equidistribution corollaries remain
\statusnew
\item \textbf{Unconditional gap equivalences.}  The conjecture is equivalent to the existence
of any nonempty open subinterval of $(0,1)$ free of fractional parts of nonintegral
solutions, and to fractional parts being bounded away from $0$, or from $1$.  \statuslean
\item \textbf{Sharper local combinatorics.}  A progression $n,n+h,\dots,n+rh$ of candidates
is impossible whenever $|\thetaq(\thetaq-1)\cdots(\thetaq-r+1)|\,h^r<n^{r-\thetaq}$,
strictly improving the committed three-term criterion $h^5\le n$ and extending it to all
lengths.  The three-term case at exponent $83/400=0.2075$ is kernel-verified
(\lean{SharpProgression}); the arbitrary-length family remains \statusnew
\item \textbf{The irreducible remainder.}  The conjecture is equivalent to: no odd integer
$u>1$ satisfies $u^{\thetaq}\in\Zthree$.  The equivalence is kernel-verified
(\lean{Localization}); the statement itself is open even for $u=3$.
\end{enumerate}

\section{The problem and its equivalent forms}
\label{sec:problem}

\subsection{Original statement}

\begin{conjecture}[Alaoglu--Erd\H{o}s, two-base integer form]
\label{conj:AE}
For every $x\in\R$,
\[
   2^x\in\Z\quad\text{and}\quad 3^x\in\Z
   \qquad\Longrightarrow\qquad x\in\Z.
\]
\end{conjecture}

\begin{statusbox}{recorded in Lean as the proposition \lean{AlaogluErdosConjecture}, without
being asserted.}
\end{statusbox}

Since $2^x$ and $3^x$ are positive, the hypotheses actually put them in $\N$; and $2^x\in\N$
forces $x\ge0$, since a negative exponent gives $0<2^x<1$.  All solutions therefore live in
$[0,\infty)$ (\lean{nonneg_of_two_rpow_integer}).

Questions of this type go back to the 1944 study of colossally abundant numbers by Alaoglu
and Erd\H{o}s \cite{AE1944}, where the arithmetic nature of simultaneous rational powers
$p^x,q^x$ of distinct primes controls the quotients of consecutive colossally abundant
numbers --- a theme reaching back to Ramanujan's superior highly composite numbers
\cite{Ram1915}.  The modern two-base integer form is elementary to state and is a direct
special case of the four exponentials conjecture (Schneider's first problem
\cite{Schneider1957}; see Waldschmidt \cite{Waldschmidt2000,Waldschmidt2023}).  As of the
dated literature check for this report, the four exponentials conjecture remains open.

\subsection{Candidate-base normalization}

Put
\[
  \thetaq:=\frac{\log 3}{\log 2}=\log_2 3
  =1.584962500721156\ldots
\]
and define the \emph{candidate set}
\[
 \Cand:=\{m\in\N_{>0}:m^{\thetaq}\in\N\}.
\]

\begin{lemma}[Exponent--candidate bijection]
\label{lem:bijection}
The maps $x\mapsto 2^x$ and $m\mapsto\log_2m$ are mutually inverse bijections between
$\Sol:=\{x\in\R:2^x,3^x\in\Z\}$ and $\Cand$.  Under the bijection, $3^x=m^{\thetaq}$.
\end{lemma}

\begin{proof}
If $x\in\Sol$, set $m=2^x\in\N$.  Then
$m^{\thetaq}=\exp(\log m\cdot\log3/\log2)=\exp(x\log3)=3^x\in\N$.  Conversely, if
$m\in\Cand$ and $x=\log_2m$, then $2^x=m$ and the same identity gives
$3^x=m^{\thetaq}\in\N$.
\end{proof}

\begin{statusbox}{\statuslean{} --- \lean{exists_twoBaseNaturalCandidate_of_two_three_rpow_integer},
\lean{TwoBaseNaturalCandidate.integer_powers}.}
\end{statusbox}

Integral solutions correspond exactly to powers of two\break
(\lean{twoBaseCandidateExponent_isInteger_iff_isPowerOfTwo}), so the conjecture is
equivalent to
\[
  \Cand=\{2^n:n\in\Nzero\}
\]
(\lean{alaogluErdosConjecture_iff_candidates_are_powers_of_two}).  The candidate formulation
is well suited to unique factorization, finite certification, finite differences of
$t^{\thetaq}$, and multiplicative-rank arguments, and it is the formulation in which most of
the corpus is stated.

Two closure properties of $\Sol$ are elementary but essential: $\Sol$ is closed under
addition and under multiplication by positive naturals, so it is an additive submonoid of
$[0,\infty)$ containing $\N$ (\lean{rpow_integer_add}, \lean{rpow_integer_nat_mul}).

\subsection{The four-exponentials reduction}

Suppose $x$ is nonintegral and $2^x,3^x$ are algebraic.  The rational-exponent argument
(Lemma~\ref{lem:rational}) shows that $x$ cannot be rational, so $1,x$ are $\Q$-linearly
independent; and $\log2,\log3$ are $\Q$-linearly independent since $2^a3^b=1$ forces
$a=b=0$.  The four exponentials built from these two pairs are
\[
 e^{1\cdot\log2}=2,\quad e^{1\cdot\log3}=3,\quad e^{x\log2}=2^x,\quad e^{x\log3}=3^x,
\]
all four algebraic --- contrary to the four exponentials conjecture.  This reduction is
formalized: \lean{alaogluErdosConjecture_of_realFourExponentialsConjecture}.  \statuslean

The reduction is a diagnostic, not just context: any proposed elementary argument must
exploit integrality, positivity, or arithmetic structure beyond bare algebraicity of those
four values, or else it would prove a substantial case of the four exponentials conjecture
by another name.  By contrast, three multiplicatively independent bases fall to the six
exponentials mechanism, and the special case needed here is fully formalized:

\begin{theorem}[Three-base theorem]
\label{thm:threebase}
If $2^x,3^x,5^x\in\Z$ for a real $x$, then $x\in\Z$.
\end{theorem}

\begin{statusbox}{\statuslean{} --- modules \lean{IntegerExponent},
\lean{IntegerExponent.DeterminantBound}, \lean{IntegerExponent.ExponentialKernel},
\lean{IntegerExponent.Grid}, following Waldschmidt's square-determinant proof
\cite{WaldschmidtSix}.}
\end{statusbox}

In outline, for irrational $x$ one forms a large square matrix
$\mathcal A_{ij}=\exp(a_ib_j)$, where row arguments encode monomials in $2,3,5$ and column
arguments encode affine forms in $1,x$.  Irrationality makes the column parameters
pairwise distinct, multiplicative independence the row parameters; a total-positivity
argument gives $\det\mathcal A\ne0$; integrality of all entries forces
$|\det\mathcal A|\ge1$; and an interpolation estimate with explicit size choices gives
$|\det\mathcal A|<1$, a contradiction.  The dimension count is critical: with three
independent row directions and two column directions the analytic bound beats the arithmetic
bound, while in the $2\times2$ case the same balance lands exactly on the classical boundary
where decay below one fails.  A successful two-base determinant would be a new
four-exponentials argument.

\subsection{Rational and algebraic exponents}

\begin{lemma}[Rational exponent]
\label{lem:rational}
If $x\in\Q$ and $2^x\in\Z$, then $x\in\Z$.  Consequently every nonintegral solution is
irrational.
\end{lemma}

\begin{proof}
Write $x=a/b$ in lowest terms with $b>0$.  If $2^{a/b}=m\in\N$, then $2^a=m^b$; comparing
the exponent of $2$ gives $a=b\,v_2(m)$, so $b\mid a$ and $b=1$.
\end{proof}

\begin{statusbox}{\statuslean{} --- \lean{IntegerExponent.integer_of_rational_of_two_rpow_integer},
\lean{irrational_of_not_integer_of_two_rpow_integer}.}
\end{statusbox}

Gelfond--Schneider upgrades irrational to transcendental; see \cref{sec:transcendence}.

\section{Repository timeline and verification status}
\label{sec:commits}

The formal development advanced in three waves plus documentation.  All commits are in the
repository \cite{ProveItRepo}; the mathematical role of each is summarized below.

\begin{longtable}{@{}p{0.13\textwidth}p{0.80\textwidth}@{}}
\toprule
Commit & Role \\
\midrule
\endfirsthead
\toprule
Commit & Role \\
\midrule
\endhead
\commit{f492c18f1} & Three-base theorem via interpolation determinants
(\lean{IntegerExponent} and submodules). \\
\addlinespace
\commit{5ed4e6185} & Merge of the three-base branch. \\
\addlinespace
\commit{a37d97f9d} & Basic two-base module: the conjecture as a proposition, closure under
natural multiples, irrationality of nonintegral solutions, endpoint gaps, numerical bounds
for $\thetaq$, the $2^x<10$ exclusion, second-difference obstruction, aligned-block $2/3$
counting. \\
\addlinespace
\commit{4a9b03d12} & Major corpus: finite checks through $2^x<256$; arithmetic descent;
multiplicative rank; candidate structure; zero density; shrinking targets; the $3$-smooth
theorem; the four-exponentials reduction; transcendence consequences; and a complete
Gelfond--Schneider port. \\
\addlinespace
\commit{18a85e6b8} & Merge of the two-base branch. \\
\addlinespace
\commit{423ac24f0} & First survey document (\lean{two-base-partial-results.tex}) with
rendered PDF. \\
\addlinespace
\commit{98646b455} & Three further modules, kernel-verified:
\lean{FractionalPartDensity} (density dichotomy and gap equivalences),
\lean{OddCore} (common odd core, improved counting, dyadic blocks),
\lean{FiniteCheck4096} (zero-free region to $2^x<4096$). \\
\addlinespace
\commit{716617d72} & Survey finalized after kernel acceptance of the $4096$ region. \\
\addlinespace
\commit{11ac3093f} & Independent research report (audit at \commit{423ac24f0}, exact
conditional structure, interval scan to $2^{20}$); one of the two sources merged into the
present document. \\
\bottomrule
\end{longtable}

The audit performed in \commit{11ac3093f} treated results of the then-uncommitted modules as
article-level claims; with \commit{98646b455} in the tree and kernel-verified, those results
are now cited as committed Lean.  The interval scan to $2^{20}$ remains at the
\statuscomp{} level by design.
\section{The kernel-verified corpus}
\label{sec:corpus}

Throughout this section $x$ denotes a hypothetical nonintegral solution, $m=2^x\in\N$,
$n=3^x\in\N$.  Every statement here is \statuslean{} unless labeled otherwise.

\subsection{Fractional-part gaps and shrinking targets}
\label{sec:gaps}

An integral power caught strictly between consecutive integral powers of its base must clear
both by at least one full unit; dividing through normalizes this to the fractional part.

\begin{theorem}[Fractional-part gaps]
\label{thm:gaps}
Let $x\in(N,N+1)$ with $N\in\N$, $t=x-N$, and suppose $2^x,3^x\in\Z$.  Then
\[
1+2^{-N}\le 2^{t}\le 2-2^{-N},
\qquad
1+3^{-N}\le 3^{t}\le 3-3^{-N}.
\]
\end{theorem}

\begin{statusbox}{\statuslean{} --- \lean{two_three_fractional_part_gaps},
\lean{two_three_rpow_gaps_near_nat}, from
\lean{rpow_integer_gap_between_consecutive_powers}.}
\end{statusbox}

Near the endpoints the excluded zones have width $\asymp2^{-N}$: they \emph{shrink} as the
integer part grows.  Applied to all multiples $kx$ (which are again solutions), this yields
explicit Diophantine lower bounds.

\begin{theorem}[Effective irrationality bound]
\label{thm:shrinking}
Let $x\notin\Z$ satisfy $2^x\in\Z$ and write $m=2^x$.  Then for all integers $p$ and all
$k\ge1$,
\[
\Bigl|x-\frac pk\Bigr|\ \ge\ \frac1{2\,k\,m^{k}} .
\]
Exact (slightly stronger) radii, and two-base versions taking the better of the base-$2$ and
base-$3$ radii, are formalized as well.
\end{theorem}

\begin{statusbox}{\statuslean{} --- \lean{rational_approximation_exponential_lower_bound},
\lean{two_three_multiple_distance_to_integer_lower_bound},
\lean{exists_natural_base_rational_approximation_lower_bound} in \lean{ShrinkingTarget}.}
\end{statusbox}

In shorthand, $\distZ{kx}\gtrsim2^{-kx}$.  This bound is compatible with ordinary Diophantine
approximation: Dirichlet guarantees infinitely many $k$ with $\distZ{kx}<1/k$, and the
shrinking radius $2^{-kx}$ is incomparably smaller.  The interplay between this exponential
gap and equidistribution is discussed in \cref{sec:attempts}.

\subsection{Verified zero-free regions and certificate architecture}
\label{sec:finite}

\begin{theorem}[Zero-free region]
\label{thm:finite4096}
If $2^x,3^x\in\Z$ and $2^x<4096$, then $x\in\Z$.  Equivalently, every nonintegral solution
satisfies $x\ge12$, and every candidate that is not a power of two is at least $4096$.
\end{theorem}

\begin{statusbox}{\statuslean{} ---
\lean{integer_of_two_three_rpow_integer_of_two_rpow_lt_4096},
\lean{twelve_le_of_not_integer_of_two_three_rpow_integer},
\lean{le_of_twoBaseNonintegerCandidate_4096} in \lean{FiniteCheck4096}, extending
\lean{FiniteCheck} ($2^x<100$) and \lean{FiniteCheck256} ($2^x<256$, hence $x\ge8$).}
\end{statusbox}

The certificate architecture is worth recording, both because it is the repository's model
of a trustworthy large computation and because scaling it is a named research direction.
For each non-power-of-two $m$ the certificate is a pair of exact integer power comparisons
against rational enclosures $L_p/L_q<\thetaq<U_p/U_q$ drawn from continued-fraction
convergents of $\thetaq$: from
\[
a^{L_q}<m^{L_p}
\quad\text{and}\quad
m^{U_p}<(a+1)^{U_q}
\]
one traps $a<m^{\thetaq}<a+1$, so $m^{\thetaq}\notin\Z$.  The $2^x<256$ stage verifies
per-value certificates by \lean{norm_num} with three enclosure tiers (exponent sizes
${\sim}500$, ${\sim}1500$, ${\sim}25000$).  The extension to $4096$ verifies $3836$
certificates through a single kernel-replayed \lean{decide} over a table generated by
PARI/GP and revalidated independently with exact big-integer arithmetic in Python.  Five
tiers suffice, with tier populations $78$, $488$, $3241$, $28$, $1$:

\begin{center}
\small
\begin{tabular}{@{}crr@{}}
\toprule
tier & lower convergent $L_p/L_q$ & upper convergent $U_p/U_q$ \\
\midrule
0 & $569/359$ & $485/306$ \\
1 & $1054/665$ & $1539/971$ \\
2 & $25781/16266$ & $24727/15601$ \\
3 & $50508/31867$ & $125743/79335$ \\
4 & $176251/111202$ & $301994/190537$ \\
\bottomrule
\end{tabular}
\end{center}

Only the value $m=2762$, whose $\thetaq$-power is within $3\cdot10^{-9}$ of an integer in
logarithmic scale, needs tier $4$; its certificate compares integers of roughly $90{,}000$
digits, which the kernel's exact arithmetic replays without difficulty.  No axioms beyond
the Lean kernel are involved; in particular no \lean{native_decide}.

\subsection{Arithmetic descent}
\label{sec:descent}

\begin{theorem}[Affine descent]
\label{thm:descent}
If the two integral outputs factor simultaneously as $2^x=2^ru^k$ and $3^x=3^rv^k$ with
$r\ge0$, $k\ge1$, then $(x-r)/k$ is again a solution.  Combined with the $x\ge8$ region this
gives the unconditional constraint
\[
r+8k\le x
\]
for nonintegral $x$; in particular common perfect-power shapes of the two outputs are
strongly restricted.
\end{theorem}

\begin{statusbox}{\statuslean{} --- \lean{affine_descent_integral_powers},
\lean{base_depth_add_eight_mul_degree_le},
\lean{min_output_base_factorization_add_eight_le} in \lean{ArithmeticRigidity}.  With the
$4096$ region the constant $8$ improves to $12$ by bookkeeping; with the interval scan of
\cref{sec:finite-computation} it would improve to $20$ at the \statuscomp{} level.}
\end{statusbox}

For the \emph{least} counterexample the same descent later becomes an exact primitivity
statement (Corollary~\ref{cor:affine-primitive}): no nontrivial matched decomposition exists at all.

\subsection{Local progression obstructions and zero density}
\label{sec:ap}

The function $t\mapsto t^{\thetaq}$ is strictly convex with second differences at moderate
steps lying strictly between $0$ and $1$ --- hence never integers.

\begin{theorem}[Committed AP obstruction]
\label{thm:ap}
Let $n,h\ge1$ with $h^5\le n$.  Then $n,n+h,n+2h$ are not all candidates.  In particular no
three consecutive naturals are candidates, and among the positive integers up to $3N$ at
most $2N$ are candidates.
\end{theorem}

\begin{statusbox}{\statuslean{} ---
\lean{not_three_arithmetic_progression_twoBaseNaturalCandidates},
\lean{not_three_consecutive_logThreeDivLogTwo_rpow_integers},
\lean{twoBaseNaturalCandidate_Icc_card_le}.  The sharper criterion
$\thetaq(\thetaq-1)h^2<n^{2-\thetaq}$ and its arbitrary-length generalization are proved in
\cref{sec:higher-differences} at the \statusnew{} level.}
\end{statusbox}

The hypothesis bounding $h$ is essential: for steps comparable to $n$ the second difference
is large and the obstruction vanishes.  This locality is a genuine barrier --- a
\emph{global} progression obstruction would, for example, forbid candidates with odd part
$2^s\pm1$ via progressions such as $(2^i,\,2^{i+t-1},\,2^i(2^t-1))$, but no such global
statement is available; under a counterexample the exact monoid of
\cref{sec:exact-structure} shows sparse candidates whose additive relations the curvature
argument cannot see.

Feeding the local obstruction into the theory of Roth numbers gives quantitative sparsity.

\begin{theorem}[Zero density]
\label{thm:density}
For every $H\ge1$ the number of candidates below $N$ is at most
$H^5+\lceil(N-H^5)/H\rceil\,r_3(H)$, where $r_3(H)$ is the maximal size of a $3$-AP-free
subset of $\{0,\dots,H-1\}$.  Consequently the candidate counting function is $o(N)$: the
candidates, and a fortiori the values $2^x$ of nonintegral solutions, have asymptotic
density zero.
\end{theorem}

\begin{statusbox}{\statuslean{} --- \lean{twoBaseNaturalCandidatesBelow_card_le},
\lean{twoBaseNaturalCandidatesBelow_isLittleO_id},
\lean{tendsto_twoBaseNonintegerCandidatesBelow_density_zero} in \lean{ZeroDensity} and
\lean{NonintegerDensity}.}
\end{statusbox}

\subsection{Multiplicative rank two and the common odd core}
\label{sec:rank}

\begin{theorem}[Rank two]
\label{thm:rank2}
The prime-exponent vectors $\mathbf v(a)=(v_p(a))_p$ of all candidates span a rational space
of dimension at most two.  Consequently at most $(\clog_2N)^2$ candidates lie below $N$.
\end{theorem}

\begin{statusbox}{\statuslean{} --- \lean{finrank_span_primeExponentVectors_le_two},
\lean{twoBaseNaturalCandidatesBelow_card_le_clog_sq} in \lean{MultiplicativeRank}.  The
mechanism abstracts the determinant method: three candidates with independent exponent
vectors would give three multiplicatively independent bases, contradicting
Theorem~\ref{thm:threebase} through the irrationality of $\thetaq$
(\lean{irrational_logThreeDivLogTwo}).}
\end{statusbox}

Since $2$ is itself a candidate, all odd parts lie on at most one further rational
direction.  Pairwise this yields the cross-valuation identity, and globally a common odd
core.

\begin{lemma}[Cross-valuation identity]
\label{lem:cross-input}
Let $a,b\in\Cand$ and let $p$ be an odd prime with $\vpp{p}{a}>0$.  Then for every odd prime
$q$,
\[
\vpp{p}{a}\,\vpp{q}{b}=\vpp{p}{b}\,\vpp{q}{a}.
\]
\end{lemma}

\begin{statusbox}{\statuslean{} --- \lean{odd_factorization_cross_mul} in
\lean{CandidateStructure}; the fixed-power relations
\lean{fixed_power_relation_of_odd_prime_factor} follow.}
\end{statusbox}

\begin{theorem}[Common odd core]
\label{thm:oddcore-lean}
There is a single odd $w>1$ such that every candidate equals $2^iw^j$ for some
$i,j\in\Nzero$; candidates that are not powers of two have $j\ge1$.  In exponent
coordinates: every solution has the form $x=i+j\log_2w$ with $i,j\in\Nzero$.
\end{theorem}

\begin{statusbox}{\statuslean{} --- \lean{exists_common_odd_core},
\lean{exists_common_odd_core_split}, \lean{exists_common_odd_core_exponent} in
\lean{OddCore}.  The canonical (gcd-normalized) version with uniqueness of $w$ and of the
representation is proved in Theorem~\ref{thm:odd-core} at the \statusnew{} level; formalizing that
normalization is Priority~A of \cref{sec:program}.}
\end{statusbox}

\begin{corollary}[Improved counting]
\label{cor:counting}
The number of candidates below $N$ is at most $\clog_2N\cdot\clog_3N$, and each dyadic
block $[2^T,2^{T+1})$ contains at most $\clog_3(2^{T+1})\approx0.631\,(T+1)$ candidates.
In exponent coordinates, the number of solutions with $x\in[T,T+1)$ grows at most linearly
in $T$, against the quadratic global bound.
\end{corollary}

\begin{statusbox}{\statuslean{} --- \lean{twoBaseNaturalCandidatesBelow_card_le_clog_mul_clog},
\lean{twoBaseNonintegerCandidatesBelow_card_le_clog_mul_clog},
\lean{twoBaseNaturalCandidatesInDyadicBlock_card_le} in \lean{OddCore}.}
\end{statusbox}

The logarithmic-square order cannot be improved by counting alone: if a candidate with an
odd prime factor exists, multiplicative closure generates injective boxes $2^ia^j$, so the
count below $2^ka^l$ is at least $kl$
(\lean{twoBaseNaturalCandidatesBelow_card_ge_box},
\lean{twoBaseNonintegerCandidatesBelow_card_ge_box}) --- a conditional lower bound of the
same quadratic-in-logarithm shape, made exact in \cref{sec:counting}.

\subsection{Transcendence}
\label{sec:transcendence}

The corpus contains a complete Lean proof of the Gelfond--Schneider theorem (Hilbert's
seventh problem), following Hua's exposition \cite{Hua1982}, ported from Michail
Karatarakis's Mathlib formalization (Mathlib PR \#42911, Apache~2.0).

\begin{theorem}[Gelfond--Schneider]
\label{thm:gs}
If $\alpha,\beta$ are algebraic, $\alpha\ne0,1$, and $\beta$ is irrational, then
$\alpha^\beta$ is transcendental.
\end{theorem}

\begin{statusbox}{\statuslean{} ---
\lean{GelfondSchneider.transcendental_cpow_of_isAlgebraic_of_irrational}.}
\end{statusbox}

\begin{corollary}[Algebraic--integral dichotomy]
\label{cor:transcendental}
If $2^x\in\Z$ then $x$ is either an integer or transcendental; algebraicity of such $x$ is
equivalent to integrality.  Every nonintegral two-base solution is transcendental, and
$\thetaq=\log_23$ is transcendental.
\end{corollary}

\begin{statusbox}{\statuslean{} --- \lean{transcendental_of_not_integer_of_two_rpow_integer},
\lean{isAlgebraic_iff_integer_of_two_rpow_integer},
\lean{transcendental_of_not_integer_of_two_three_rpow_integer},
\lean{transcendental_logThreeDivLogTwo}.}
\end{statusbox}

The dichotomy does not resolve the conjecture --- the unknown exponent is allowed to be
transcendental --- but it constrains every construction below: the primitive generator
$\beta$ of \cref{sec:exact-structure} is transcendental, as is $\log_2w$ for the odd core.

\subsection{Classification of auxiliary bases: the \texorpdfstring{$3$}{3}-smooth theorem}
\label{sec:threesmooth}

\begin{theorem}[$3$-smooth classification]
\label{thm:threesmooth}
Suppose $x\notin\Z$ and $2^x,3^x\in\Z$.  Then for a positive natural $a$, the power $a^x$ is
an integer if and only if $a=2^u3^v$ is $3$-smooth.  In particular a third base with any
prime factor other than $2$ and $3$ forces $x\in\Z$.
\end{theorem}

\begin{statusbox}{\statuslean{} ---
\lean{rpow_integer_iff_eq_two_pow_mul_three_pow_of_not_integer},
\lean{integer_of_two_three_a_rpow_integer_of_prime_factor} in \lean{ThreeSmooth}.}
\end{statusbox}

This is an exact classification of bases for a \emph{fixed} hypothetical counterexample.  It
should not be confused with the classification of candidate values $2^x$ as $x$ ranges over
all solutions, which is the subject of the next section.  Note also the scope limitation
explained in \cref{sec:attempts}: the theorem controls the quantity $b^x$, not divisibility
of the integer $2^x$ by $b$.

\subsection{Unconditional gap equivalences}
\label{sec:dichotomy}

Because natural multiples of a solution are solutions, one nonintegral solution generates
infinitely many, with irrational generator; Dirichlet approximation and a rotation-stepping
argument spread their fractional parts densely.

\begin{theorem}[Density dichotomy]
\label{thm:dichotomy}
If some $x\notin\Z$ satisfies $2^x,3^x\in\Z$, then for every $t\in[0,1]$ and
$\varepsilon>0$ there is a nonintegral solution $y$ with $|\fract{y}-t|<\varepsilon$; and
the set of nonintegral solutions is infinite.
\end{theorem}

\begin{theorem}[Gap equivalences]
\label{thm:equiv}
Each of the following is equivalent to Conjecture~\ref{conj:AE}:
\begin{enumerate}[label=\textup{(\roman*)}]
\item some nonempty open interval $(u,v)\subseteq(0,1)$ contains no fractional part of a
nonintegral solution;
\item the fractional parts of nonintegral solutions are bounded away from $1$;
\item the fractional parts of nonintegral solutions are bounded away from $0$.
\end{enumerate}
\end{theorem}

\begin{statusbox}{\statuslean{} --- \lean{dense_fract_of_noninteger_solution},
\lean{infinite_noninteger_solutions_of_exists},
\lean{alaogluErdosConjecture_iff_fract_avoids_interval},
\lean{alaogluErdosConjecture_iff_fract_bounded_away_one},
\lean{alaogluErdosConjecture_iff_fract_bounded_away_zero} in
\lean{FractionalPartDensity}, built on the reusable
\lean{exists_pos_nat_fract_close_of_irrational}.}
\end{statusbox}

\begin{remark}[Consistency pincer]
\label{rem:pincer}
Theorem~\ref{thm:gaps} bounds $\fract{x}$ away from $0$ and $1$ by margins
$\asymp2^{-\lfloor x\rfloor}$ that decay with the height of the solution, while
Theorem~\ref{thm:dichotomy} produces fractional parts arbitrarily close to $0$ and $1$ along
multiples of unbounded height.  The two results meet edge-to-edge without contradiction, and
Theorem~\ref{thm:equiv} shows the decay is essential: \emph{any} height-independent gap, however
small and wherever located, is exactly as strong as the full conjecture.  The elementary gap
bounds are therefore of optimal shape, and the fractional-part axis measures the entire
remaining distance to Conjecture~\ref{conj:AE}.
\end{remark}

These equivalences are \emph{unconditional} biconditionals.  Conditional on failure, the
exact structure of the next section strengthens density to blockwise Weyl equidistribution
(Theorem~\ref{thm:block-equidist}).

\section{Exact conditional structure of all solutions}
\label{sec:exact-structure}

This section develops the exact rigidity theory of a
hypothetical counterexample.  It starts from the kernel-verified cross-valuation identity
(Lemma~\ref{lem:cross-input}) and sharpens the verified odd-core theorem
(Theorem~\ref{thm:oddcore-lean}) in two steps: a canonical (primitive) normalization of the core,
and the removal of all residual freedom in the exponent pairs.

\subsection{The primitive common odd core}

\begin{theorem}[Primitive common odd core]
\label{thm:odd-core}
Assume $\Cand$ contains a non-power of two.  Then there is a unique odd integer $w>1$ that
is power-free in the sense
\[
\gcd\{\vpp{p}{w}:p\mid w\}=1,
\]
and every candidate $b\in\Cand$ has a unique representation
\[
b=2^iw^j,\qquad i,j\in\Nzero.
\]
\end{theorem}

\begin{proof}
Choose a non-power-of-two candidate $a$.  Its odd valuation vector
$A=(A_q)_{q\ne2}$, $A_q=\vpp qa$, is nonzero and finitely supported.  Let
\[
g=\gcd\{A_q:A_q>0\},\qquad c_q=A_q/g,\qquad w=\prod_{q\ne2}q^{c_q}.
\]
Then $w$ is odd, greater than one, and the positive $c_q$ have gcd one.

Let $b\in\Cand$ and set $B_q=\vpp qb$.  Fix an odd prime $p$ with $A_p>0$.  By
Lemma~\ref{lem:cross-input}, $A_pB_q=B_pA_q$ for all odd $q$.  If $A_q=0$, this forces $B_q=0$.
On the support of $A$ it says
\[
B_q=\frac{B_p}{A_p}A_q=\Bigl(\frac{gB_p}{A_p}\Bigr)c_q .
\]
Put $t=gB_p/A_p\in\Q_{\ge0}$.  Every $tc_q$ is an integer.  Since the $c_q$ have gcd one,
there are integers $r_q$, finitely many nonzero, with $\sum r_qc_q=1$; therefore
$t=\sum_qr_q(tc_q)\in\Z$.  Writing $j=t\in\Nzero$, the odd part of $b$ is exactly $w^j$, so
$b=2^iw^j$ with $i=v_2(b)$.

Uniqueness of $(i,j)$ for fixed $w$ follows from odd valuations and the $2$-adic valuation.
If $w'$ is another primitive common core, the odd valuation vectors of $w$ and $w'$ generate
the same rational ray and are both primitive integer vectors, hence equal.
\end{proof}

The gcd normalization is what makes $w$ canonical: it prevents hidden perfect powers from
being absorbed into the exponent.  The kernel-verified \lean{exists_common_odd_core}
provides the existence of \emph{some} common core; the primitive normalization and
uniqueness above are the part awaiting formalization.

Set
\[
\alpha:=\log_2w,\qquad E:=3^{\alpha}=w^{\thetaq}.
\]
Every candidate exponent then has the form $x=i+j\alpha$.  The number $\alpha$ is
irrational (else the rational-exponent lemma applied to $2^\alpha=w$ would make it integral,
forcing the odd $w>1$ to be a power of two), and indeed transcendental by
Corollary~\ref{cor:transcendental}.  The candidate condition becomes
\[
3^{i+j\alpha}=3^iE^j\in\N,
\]
and at this stage the admissible pairs $(i,j)$ still form an unknown subset of $\Nzero^2$.
The next subsection shows this subset is completely rigid.

\subsection{The rational-power index and the least generator}

Because a non-power-of-two candidate exists, some $j>0$ has $E^j\in\Q$.  Let
\[
d:=\min\{j\in\N_{>0}:E^j\in\Q\}.
\]

\begin{lemma}[All rational powers are multiples of the least one]
\label{lem:rational-power-index}
For every $j\in\Z$:\quad $E^j\in\Q\iff d\mid j$.
\end{lemma}

\begin{proof}
$J:=\{j\in\Z:E^j\in\Q\}$ is an additive subgroup of $\Z$ ($0\in J$; products add exponents;
reciprocals negate them), and every nonzero subgroup of $\Z$ is $d\Z$ for its least positive
element.
\end{proof}

Write the positive rational $E^d$ in lowest terms.  Its denominator has no prime other than
$3$: taking a candidate with odd-core exponent $j=dk>0$ and $2$-exponent $i$, integrality of
$3^i(E^d)^k$ shows any other prime in the reduced denominator would survive in the $k$-th
power.  Hence there are unique integers $a\ge0$, $A\ge1$ with
\begin{equation}
E^d=\frac A{3^a},
\label{eq:Ed}
\end{equation}
and $3\nmid A$ when $a>0$.  Define
\begin{equation}
\beta:=a+d\alpha=a+d\log_2w,\qquad
M:=2^{\beta}=2^aw^d,\qquad
A=3^{\beta}=3^aE^d .
\label{eq:betaMA}
\end{equation}
Thus $\beta$ is itself a solution, irrational because $\alpha$ is.

\begin{theorem}[Exact primitive-generator theorem]
\label{thm:generator}
Assume Conjecture~\ref{conj:AE} is false.  With $w,d,a,\beta,M,A$ as above:
\begin{enumerate}[label=\textup{(\roman*)}]
\item every solution has a unique representation $x=n+k\beta$ with $n,k\in\Nzero$;
\item every candidate has a unique representation $m=2^nM^k$ with $n,k\in\Nzero$;
\item every output pair has a unique representation $(2^x,3^x)=(2^nM^k,3^nA^k)$;
\item $\beta$ is the unique least nonintegral solution.
\end{enumerate}
\end{theorem}

\begin{proof}
Let $x\in\Sol$.  By Theorem~\ref{thm:odd-core}, $2^x=2^iw^j$ for unique $i,j\in\Nzero$, so
$x=i+j\alpha$.  From $3^x=3^iE^j\in\N$ we get $E^j\in\Q$, so
Lemma~\ref{lem:rational-power-index} gives $j=dk$.  Using \eqref{eq:Ed},
\[
3^x=3^i(E^d)^k=3^{\,i-ak}A^k .
\]
Because $A/3^a$ is reduced, this is an integer exactly when $i\ge ak$.  Put
$n=i-ak\in\Nzero$; then
\[
x=i+dk\alpha=n+k(a+d\alpha)=n+k\beta .
\]
Conversely, for any $n,k\in\Nzero$,
\[
2^{\,n+k\beta}=2^nM^k\in\N,\qquad 3^{\,n+k\beta}=3^nA^k\in\N,
\]
so $n+k\beta\in\Sol$.  If $n+k\beta=n'+k'\beta$ with $k\ne k'$, then
$\beta=(n'-n)/(k-k')\in\Q$, a contradiction; hence the representation is unique.  This
proves (i), and (ii)--(iii) follow by exponentiation.  A nonintegral solution has $k\ge1$;
the least value is attained at $(n,k)=(0,1)$ and equals $\beta$, uniquely so by uniqueness
of representation.
\end{proof}

\begin{statusbox}{\statuslean{} --- \lean{exists_primitive_generator} in
\lean{PrimitiveGenerator}, formalized from the kernel-verified common odd core (the proof
does not need the canonical normalization).  The principal structural result of this
report.  It
sharpens the kernel-verified rank-two and odd-core theorems: conditional on failure, $\Sol$
is not merely contained in a rank-two additive set --- it is the free additive monoid on $1$
and one irrational generator $\beta$.  In particular the ``exponent semigroup'' degree of
freedom contemplated in the earlier survey does not exist: the set of occurring odd-core
exponents is exactly $d\,\Nzero$, with the power of two governed by one linear threshold.}
\end{statusbox}

Packaged algebraically: $(\Nzero^2,+)\xrightarrow{\ \sim\ }(\Sol,+)$ via
$(n,k)\mapsto n+k\beta$, and $(\Nzero^2,+)\xrightarrow{\ \sim\ }(\Cand,\cdot)$ via
$(n,k)\mapsto2^nM^k$; the output-pair monoid $\Pair$ is freely generated by $(2,3)$ and
$(M,A)$.  Every multiplicative relation among candidates is the obvious coordinate
relation.  This exact conditional model is the benchmark against which any future density,
progression, or descent argument should be tested.

\subsection{Primitivity of the least output pair}

Leastness of $\beta$ upgrades the repository's descent inequalities to exact exclusions.

\begin{corollary}[No matched base depth]
\label{cor:no-depth}
For the primitive pair $(M,A)=(2^{\beta},3^{\beta})$:
$\min\{v_2(M),v_3(A)\}=0$; that is, $M$ is odd or $3\nmid A$ (possibly both).
\end{corollary}

\begin{proof}
If both valuations were positive, then $2^{\beta-1}=M/2$ and $3^{\beta-1}=A/3$ would be
integers, so $\beta-1$ would be a smaller nonintegral solution.
\end{proof}

\begin{corollary}[No common perfect-power degree]
\label{cor:no-common-power}
There is no $k\ge2$ for which both $M$ and $A$ are perfect $k$-th powers.
\end{corollary}

\begin{proof}
If $M=U^k$, $A=V^k$, then $2^{\beta/k}=U$ and $3^{\beta/k}=V$, and $\beta/k$ is a smaller
nonintegral solution.
\end{proof}

\begin{corollary}[Exact affine primitivity]
\label{cor:affine-primitive}
If $M=2^rU^k$ and $A=3^rV^k$ with $r\in\Nzero$, $k\in\N_{>0}$, $U,V\in\N$, then $r=0$ and
$k=1$.
\end{corollary}

\begin{proof}
Affine descent (Theorem~\ref{thm:descent}) gives the solution $(\beta-r)/k$, nonintegral since
$\beta$ is irrational; if $r>0$ or $k>1$ it is strictly smaller than $\beta$.
\end{proof}

\subsection{Exact counting}
\label{sec:counting}

\begin{theorem}[Exact unit-block count]
\label{thm:block-count}
Assume a counterexample exists and let $\beta$ be the primitive generator.  For every
$N\in\Nzero$, the nonintegral solutions in $[N,N+1)$ are precisely
\[
x_{N,k}=\lceil N-k\beta\rceil+k\beta,
\qquad 1\le k\le\Bigl\lfloor\frac{N+1}{\beta}\Bigr\rfloor,
\]
so their number is exactly $\lfloor(N+1)/\beta\rfloor$, and
\[
\#\bigl(\Cand\cap[2^N,2^{N+1})\bigr)=1+\Bigl\lfloor\frac{N+1}{\beta}\Bigr\rfloor .
\]
\end{theorem}

\begin{proof}
For fixed $k\ge1$ there is at most one integer $n$ with $N\le n+k\beta<N+1$; such $n\ge0$
exists iff $k\beta<N+1$, and then $n=\lceil N-k\beta\rceil$.  Irrationality of $\beta$ makes
$(N+1)/\beta$ non-integral, giving the stated range; the one integral solution in the block
is $N$.
\end{proof}

With the verified bound $\beta>12$ the block count is at most $\lfloor N/12\rfloor$
(and at most $\lfloor N/20\rfloor$ under the interval scan).  Summing over blocks, the
cumulative candidate count below $2^T$ is
\begin{equation}
C(T)=T+\sum_{r=1}^{T}\Bigl\lfloor\frac r\beta\Bigr\rfloor
=\frac{T^2}{2\beta}+O(T).
\label{eq:exact-cumulative}
\end{equation}
This is conditionally sharper than the verified universal bounds ($(\clog_2N)^2$ and
$\clog_2N\cdot\clog_3N$) and matches the verified conditional lower boxes in leading order:
with $\beta>12$ the leading coefficient is below $1/24$.

\subsection{Blockwise equidistribution}

The fractional part of $x_{N,k}$ is $\fract{k\beta}$: a high unit interval contains an
initial segment of one irrational rotation orbit.

\begin{theorem}[Blockwise Weyl equidistribution]
\label{thm:block-equidist}
Assume a counterexample exists.  For every interval $I\subseteq[0,1]$ whose boundary has
measure zero,
\[
\frac{\#\{x\in(\Sol\setminus\Nzero)\cap[N,N+1):\fract x\in I\}}
{\#\bigl((\Sol\setminus\Nzero)\cap[N,N+1)\bigr)}\longrightarrow|I|
\qquad(N\to\infty).
\]
\end{theorem}

\begin{proof}
By Theorem~\ref{thm:block-count} the relevant fractional parts are
$\fract\beta,\fract{2\beta},\dots,\fract{K_N\beta}$ with
$K_N=\lfloor(N+1)/\beta\rfloor\to\infty$.  For every nonzero integer $h$ the Weyl sum is a
finite geometric series,
\[
\frac1K\sum_{k=1}^{K}e^{2\pi ihk\beta}
=\frac{e^{2\pi ih\beta}\bigl(1-e^{2\pi ihK\beta}\bigr)}{K\bigl(1-e^{2\pi ih\beta}\bigr)}
\longrightarrow0,
\]
since $\beta$ is irrational.  Weyl's criterion \cite{KuipersNiederreiter} concludes.
\end{proof}

\begin{statusbox}{\statusnew{} --- strictly stronger than the kernel-verified density
statement of Theorem~\ref{thm:dichotomy}, conditional on failure.  No effective discrepancy rate
follows without further Diophantine input on $\beta$: the geometric-sum proof controls each
Fourier mode but not uniformly over many modes.}
\end{statusbox}
\section{Equivalent one-number and radical formulations}
\label{sec:radical}

\subsection{Localization at the prime three}

\begin{theorem}[Localization reformulation]
\label{thm:localization}
Conjecture~\ref{conj:AE} is false if and only if there exists an odd integer $u>1$ with
\[
u^{\thetaq}\in\Zthree,
\]
equivalently, an odd $u>1$ and $a\in\Nzero$ with $3^au^{\thetaq}\in\N$.
\end{theorem}

\begin{proof}
If $x$ is a counterexample, write $2^x=2^au$ with $u$ odd; since $2^x$ is not a power of
two, $u>1$.  Using $(2^a)^{\thetaq}=3^a$,
\[
u^{\thetaq}=\frac{(2^x)^{\thetaq}}{(2^a)^{\thetaq}}=\frac{3^x}{3^a}\in\Zthree .
\]
Conversely, if $u^{\thetaq}=B/3^a$ with $B\in\N$, set $x=a+\log_2u$; then $2^x=2^au\in\N$
and $3^x=3^au^{\thetaq}=B\in\N$, and $x\notin\Z$ since otherwise the odd $u>1$ would be a
power of two.
\end{proof}

\begin{statusbox}{\statuslean{} ---
\lean{alaogluErdosConjecture_iff_no_odd_multiple_candidate} and
\lean{alaogluErdosConjecture_iff_odd_rpow_not_localized} in \lean{Localization}.  A compact
target for future work: prove
$u^{\log_23}\notin\Zthree$ for every odd $u>1$.  The allowance of $3$-power denominators is
exactly the freedom created by shifting the exponent by an integer.}
\end{statusbox}

\subsection{Canonical radical algebraicity}

In the notation of Theorem~\ref{thm:generator}, $E=3^{\log_2w}$ satisfies $E^d=A/3^a\in\Q$.
Minimality of $d$ has an algebraic consequence.

\begin{proposition}[Exact radical degree]
\label{prop:radical-degree}
The binomial $X^d-A/3^a$ is irreducible over $\Q$, and $E$ is algebraic of exact degree
$d$.
\end{proposition}

\begin{proof}
For a positive rational $q$, the binomial $X^d-q$ can be reducible only if $q$ is a $p$-th
power in $\Q$ for some prime $p\mid d$ (the additional $4\mid d$ exceptional case concerns
$-4$ times a fourth power and is irrelevant for $q>0$).  If $q=r^p$ with $p\mid d$, the
positive real $E^{d/p}$ is the positive $p$-th root of $q$, hence rational, contradicting
minimality of $d$.
\end{proof}

A counterexample therefore produces the very special algebraic number $E=3^{\log_2w}$,
reached from a transcendental exponent.  Taking logarithms of $E^d=A/3^a$:
\begin{equation}
d\,\log w\,\log3=(\log A-a\log3)\log2 .
\label{eq:quadratic-logs}
\end{equation}
This is a \emph{bilinear} relation among logarithms of integers.  Baker's theorem controls
linear forms in logarithms with algebraic coefficients \cite{Baker1975}; it does not rule
out \eqref{eq:quadratic-logs}.  The four exponentials conjecture does.  The equation is a
precise description of the transcendence-theoretic wall, not a disguised elementary
factorization problem.

\subsection{Where the wall stands in the smallest case}

The smallest structural question already meets the open core.  A $3$-smooth candidate
$m=2^a3^b$ ($b\ge1$) would require $3^aE_3^{\,b}\in\N$ for
\[
E_3=3^{\log_23}=e^{(\log3)^2/\log2}=5.7045224946911176\ldots,
\]
and even the irrationality of $E_3$ is open: it is a ``three-logarithms'' quantity outside
the reach of Gelfond--Schneider, Baker's method, and the six exponentials theorem.  No
currently known technique excludes odd core $w=3$; this is why the corpus emphasizes
counting, structure, and verified finite regions rather than infinite families of
exclusions, and why Conjecture~\ref{conj:radical-target} below is the honest name of the obstruction.

\begin{conjecture}[Radical four-exponentials target]
\label{conj:radical-target}
For every odd integer $u>1$, $u^{\log_23}\notin\Zthree$.  Equivalently, for every primitive
odd $w>1$ and every $d\ge1$, $(3^{\log_2w})^d\notin\Zthree$.
\end{conjecture}

By Theorem~\ref{thm:localization} both formulations are exactly equivalent to Conjecture~\ref{conj:AE}.  A
proof of the special case $u=3$ --- even of the weaker assertion that $3^{\log_23}$ is
irrational --- by any genuinely new mechanism would already be a meaningful breakthrough.

\section{Sharper finite-difference obstructions}
\label{sec:higher-differences}

The committed local theorem (Theorem~\ref{thm:ap}) forbids three candidates in progression under
the convenient condition $h^5\le n$.  The curvature argument in fact yields a sharper
natural exponent and extends to arbitrary length.

For a function $f$ and step $h>0$ write $\Delta_hf(t)=f(t+h)-f(t)$ and
$\Delta_h^r=\Delta_h\circ\cdots\circ\Delta_h$.  Iterating the fundamental theorem of
calculus,
\begin{equation}
\Delta_h^rf(n)=\int_{[0,h]^r}f^{(r)}(n+t_1+\cdots+t_r)\,dt_1\cdots dt_r .
\label{eq:finite-diff-integral}
\end{equation}
For $f(t)=t^{\thetaq}$ and $r\ge2$,
$f^{(r)}(t)=\thetaq(\thetaq-1)\cdots(\thetaq-r+1)\,t^{\thetaq-r}$; since $1<\thetaq<2$ this
has a fixed nonzero sign on $(0,\infty)$ and its absolute value decreases in $t$.

\begin{theorem}[Higher-difference progression obstruction]
\label{thm:higher-difference}
Let $r\ge2$, $n,h\in\N_{>0}$, and $c_r=|\thetaq(\thetaq-1)\cdots(\thetaq-r+1)|$.  If all of
$n,n+h,\dots,n+rh$ are candidates, then $1\le c_rh^rn^{\thetaq-r}$.  Such a progression is
therefore impossible whenever
\begin{equation}
c_rh^r<n^{\,r-\thetaq}.
\label{eq:progression-condition}
\end{equation}
\end{theorem}

\begin{proof}
If all $n+jh$ are candidates, every $(n+jh)^{\thetaq}$ is an integer, so the integer
combination $\Delta_h^rf(n)$ is an integer; by \eqref{eq:finite-diff-integral} and the fixed
sign of $f^{(r)}$ it is nonzero, hence of absolute value at least one.  Bounding
$|f^{(r)}|$ by its value at $n$,
\[
1\le|\Delta_h^rf(n)|\le c_rh^rn^{\thetaq-r}. \qedhere
\]
\end{proof}

\begin{statusbox}{\statusnew{} --- an attractive Lean target: the iterated-integral identity
replaces nested mean-value constructions, and the numerical side reduces to monotonicity of
a negative real power.}
\end{statusbox}

For $r=2$ this gives the sharp three-term criterion
\begin{equation}
\thetaq(\thetaq-1)h^2<n^{2-\thetaq},
\qquad\text{i.e.}\qquad
h<1.038547802\ldots\cdot n^{0.2075187\ldots},
\label{eq:sharp-3ap}
\end{equation}
strictly weaker as a hypothesis than the committed $h\le n^{1/5}$.  The first thresholds
$h<C_rn^{1-\thetaq/r}$, $C_r=c_r^{-1/r}$:

\begin{center}
\begin{tabular}{@{}rrrr@{}}
\toprule
$r$ & length $r+1$ & exponent $1-\thetaq/r$ & $C_r$ \\
\midrule
2 & 3 & 0.207518750 & 1.038547802 \\
3 & 4 & 0.471679166 & 1.374849673 \\
4 & 5 & 0.603759375 & 1.164124485 \\
5 & 6 & 0.683007500 & 0.946705202 \\
6 & 7 & 0.735839583 & 0.778537835 \\
7 & 8 & 0.773576786 & 0.652646038 \\
8 & 9 & 0.801879687 & 0.557368999 \\
\bottomrule
\end{tabular}
\end{center}

As $r$ grows the admissible step exponent approaches one.  This looks tantalizing but does
not contradict the exact candidate monoid: in a dyadic block near $X$ the conditional number
of candidates is $O(\log X)$, so the average additive gap is of order $X/\log X$, which for
each fixed $r$ exceeds $X^{1-\thetaq/r}$.  Long-progression leverage would need $r$ growing
with the height together with control of the rapidly growing $c_r$ and a combinatorial
source of such progressions.

\section{Independent finite verification through \texorpdfstring{$2^{20}$}{2\^{}20}}
\label{sec:finite-computation}

\begin{proposition}[Finite exclusion under interval-software trust]
\label{prop:finite20}
For every integer $1\le m<2^{20}$: if $m^{\thetaq}\in\Z$ then $m$ is a power of two.
Equivalently, $2^x,3^x\in\Z$ and $2^x<2^{20}$ imply $x\in\Z$; a nonintegral solution
satisfies $x>20$.
\end{proposition}

\begin{statusbox}{\statuscomp{} --- CPython 3.13.5, \lean{mpmath} 1.3.0 interval arithmetic
at 40 decimal digits, outward endpoints.  Reproducible but \emph{not} replayed by the Lean
kernel; not interchangeable with
\lean{twelve_le_of_not_integer_of_two_three_rpow_integer}.}
\end{statusbox}

For each non-power of two $m$, a floating approximation proposes
$n=\lfloor m^{\thetaq}\rfloor$; the proposal is then \emph{certified} by interval proofs of
the two strict inequalities
\begin{align}
\log m\,\log3-\log n\,\log2&>0,\label{eq:lower-log}\\
\log(n+1)\,\log2-\log m\,\log3&>0,\label{eq:upper-log}
\end{align}
which together say exactly $n<m^{\thetaq}<n+1$.  Once both interval lower endpoints are
positive the floating proposal has no remaining logical role.  All $1{,}048{,}555$
non-powers of two below $2^{20}$ passed, in two ranges:

\begin{center}
\small
\begin{tabularx}{\textwidth}{@{}lrrX@{}}
\toprule
Range & Checked & Failures & Smallest certified log margins \\
\midrule
$1\le m<2^{19}$ & 524{,}268 & 0 &
lower $1.4147\times10^{-15}$ at $(m,n)=(189427,231498127)$;\ \
upper $3.0841\times10^{-15}$ at $(306982,497586055)$. \\
\addlinespace
$2^{19}\le m<2^{20}$ & 524{,}287 & 0 &
lower $9.4902\times10^{-16}$ at $(958460,3023921230)$;\ \
upper $2.3020\times10^{-15}$ at $(556061,1275861767)$. \\
\bottomrule
\end{tabularx}
\end{center}

At the globally tightest case the 40-digit interval had width about $1.3\times10^{-54}$,
some thirty-nine orders of magnitude below the certified positive margin; a separate
100-digit computation gives
$958460^{\thetaq}-3023921230=4.14020498680340526\times10^{-6}$.  The verification script,
full outputs, and SHA-256 manifests are reproduced in Appendix~\ref{app:interval-output}.

\paragraph{From scan to kernel certificate.}
The repository's own architecture (\cref{sec:finite}) shows the stronger trust model:
generate rational certificates externally, revalidate them independently, and replay only
exact integer comparisons in Lean by \lean{decide}.  A scalable $2^{20}$ formal certificate
should share convergent enclosures across ranges, chunk the data modules to bound
elaboration memory, and rely on verified binary powering rather than giant literals; the
worst cases identified by the scan (the four extremal pairs above) predict the required tier
depth.

\section{Attempts at a full proof and their exact failure points}
\label{sec:attempts}

\subsection{Collapse rank two to rank one}
The rank argument leaves a perfectly coherent rank-two possibility, now known exactly:
$\Cand=\langle2,M\rangle$.  Nothing in the rank machinery distinguishes this from the
allowed monoid of powers of two; excluding the second generator \emph{is} the original
transcendence problem.

\subsection{Density of multiples against endpoint gaps}
The orbit $k\beta\bmod1$ is equidistributed (Theorem~\ref{thm:block-equidist}), but the
integrality gap at height $k$ is of exponential scale $M^{-k}$
(Theorem~\ref{thm:shrinking}).  Equidistribution controls fixed intervals --- and with effective
discrepancy, polynomially shrinking ones along subsequences --- but numbers with bounded
partial quotients show that equidistributed multiples can avoid exponentially small targets
forever.  Closing this route needs Diophantine information specific to the transcendental
$\beta$ with $2^{\beta},3^{\beta}\in\N$, not generic metric theory.

\subsection{Exact counts against local progression obstructions}
Conditional failure gives $O(\log X)$ candidates per multiplicative block --- far below the
positive-density regime of Roth-type theorems --- while
Theorem~\ref{thm:higher-difference} reaches steps $X^{1-\thetaq/r}$ against average gaps
$X/\log X$.  Sparse exact structure and local curvature are mutually consistent.  A
three-term progression of candidates would solve the $S$-unit equation
$2^{n_1}M^{k_1}+2^{n_3}M^{k_3}=2^{n_2+1}M^{k_2}$; general $S$-unit theory yields finiteness
statements, not the required nonexistence.

\subsection{Exploit the primitive output pair}
The pair $(M,A)$ admits no matched depth, no common perfect power, and no nontrivial affine
decomposition (Corollaries~\ref{cor:no-depth}, \ref{cor:no-common-power} and~\ref{cor:affine-primitive}).  What remains
is exactly relation \eqref{eq:quadratic-logs}, a product of logarithms outside the scope of
linear-forms machinery.

\subsection{Force algebraicity of a forbidden auxiliary exponential}
The five exponentials theorem \cite{WaldschmidtVariations,Roy1992} makes $e^{\gamma x}$
transcendental for every nonzero algebraic $\gamma$, given the four algebraic corner
exponentials.  The natural auxiliary quantities $p^x=e^{x\log p}$ have transcendental
coefficient $\log p$, so the theorem does not directly convert divisors of $M$ or $A$ into
a contradiction.  The six exponentials theorem shows $b^x$ is transcendental for any base
$b$ multiplicatively independent of $2,3$; the committed Theorem~\ref{thm:threesmooth} is its
formalized integer-level shadow.  Neither prevents $b$ from dividing the integer $2^x$;
they control the distinct quantity $b^x$.

\section{Consolidated research program}
\label{sec:program}

The two source documents proposed overlapping programs; the exact structure theory settles
some directions and re-prioritizes the rest.  Ordered by leverage and feasibility:

\subsection*{Priority A: formalize the exact conditional structure}
\emph{Substantially completed since the first version of this report.}  The module
\lean{PrimitiveGenerator} formalizes the generator theorem itself
(\lean{exists_primitive_generator}): conditional on a nonintegral solution, it constructs
an irrational $\beta$ with $2^{\beta},3^{\beta}\in\N$ such that the solution set is exactly
$\{n+k\beta\}$ with unique representations and $\beta$ least --- using the kernel-verified
common odd core rather than the canonical normalization, which the proof turns out not to
need.  The module \lean{Localization} formalizes the localization reformulation
(\lean{alaogluErdosConjecture_iff_no_odd_multiple_candidate},
\lean{alaogluErdosConjecture_iff_odd_rpow_not_localized}), giving
Theorem~\ref{thm:localization} kernel-verified status.  Remaining from this priority:
the canonical (gcd-primitive) core with uniqueness of $w$, exact counting formulas, and
blockwise equidistribution (\lean{ExactCounting},
\lean{FractionalPartEquidistribution}).

\subsection*{Priority B: formalize the sharper finite differences}
\emph{The three-term case is now kernel-verified.}  The module \lean{SharpProgression}
proves \lean{not_three_arithmetic_progression_twoBaseNaturalCandidates_sharp}: no three
candidates in progression once $h^{400}\le n^{83}$, i.e.\ step exponent $83/400=0.2075$
against the optimal $0.20751\ldots$ of \eqref{eq:sharp-3ap} --- replacing the committed
$h^5\le n$ (exponent $0.2$).  The ingredients are the kernel-replayed enclosure
$\thetaq<317/200$ (from $3^{200}<2^{317}$) and the committed second-difference bound.
Remaining: the arbitrary-length family of Theorem~\ref{thm:higher-difference} via the
iterated-integral identity \eqref{eq:finite-diff-integral}.

\subsection*{Priority C: kernel-replay the \texorpdfstring{$2^{20}$}{2\^{}20} region}
The scan shows the finite truth is computationally routine; the bottleneck is a compact
kernel-replayable proof object.  Adaptive convergent enclosures, chunked certificate
modules, and shared tiers along ranges are the ingredients; the target theorem is the
$1048576$ analogue of Theorem~\ref{thm:finite4096}.  Every descent constant then improves to $20$
by bookkeeping.

\subsection*{Priority D: isolate the radical target in a reusable API}
Formalize Theorem~\ref{thm:localization} and Proposition~\ref{prop:radical-degree}: failure
$\iff$ an odd $u>1$ with $u^{\thetaq}\in\Zthree$ $\iff$ a primitive odd $w$, index $d$,
and reduced $q\in\Zthree$ with $(3^{\log_2w})^d=q$, plus irreducibility of $X^d-q$ and the
symmetric base-swapped version (Appendix~\ref{app:symmetric}).  This creates a sharply stated
endpoint for importing future transcendence results.

\subsection*{Priority E: five and sharp six exponentials}
Formalizing the five exponentials theorem or Roy's sharp six exponentials theorem
\cite{Roy1992} would add restrictions such as transcendence of $e^{\gamma x}$ for algebraic
$\gamma\ne0$.  The effort is comparable to the completed Gelfond--Schneider port and should
be undertaken with a specific downstream lemma in mind.

\subsection*{Priority F: auxiliary-function methods tailored to integrality}
A route beyond four exponentials must exploit more than algebraicity: here all corner
values are positive integers, all their powers remain integers, and the primitive pair is
canonical.  Promising variants: use several non-Archimedean valuations simultaneously; use
Corollary~\ref{cor:affine-primitive} to prevent determinant factors from sharing perfect powers;
work with norms in the explicit radical field $\Q(E)$ of Proposition~\ref{prop:radical-degree},
combining Archimedean smallness with denominator support only at $3$; search for a
product-formula contradiction across all places of $\Q(E)$.  The last is the most
structurally aligned with the radical formulation.

\subsection*{Priority G: primitive-pair search instead of raw candidate scanning}
A counterexample has a primitive pair $(M,A)$ with $M^{\thetaq}=A$,
$\min(v_2M,v_3A)=0$, no common perfect power, and no matched affine decomposition.  Finite
searches can reject candidates early on these conditions and record near misses keyed by
odd core, revealing whether closest approaches cluster by smoothness pattern, denominator,
or rational-power index $d$.

\subsection*{Priority H: quantitative irrationality and discrepancy}
Effective irrationality measures for $\thetaq$ (P\'ade/hypergeometric methods) would
replace per-value certificate tiers by a uniform polynomial criterion and likely extend
verified regions by orders of magnitude.  They do not directly control the generator
$\beta=\log_2M$ for unknown $M$; that would require an effective measure derived from the
relation \eqref{eq:quadratic-logs}, which lies beyond ordinary linear forms in logarithms.

\subsection*{Downstream: the colossally abundant interface}
Formalizing the original application --- the conjecture's rational-power form implies that
quotients of consecutive colossally abundant numbers are prime \cite{AE1944} --- connects
the corpus to its historical source.  It is logically downstream and should follow the core
structure modules.

\section{Conclusions}
\label{sec:conclusion}

The repository contains a broad, kernel-verified attack on the Alaoglu--Erd\H{o}s
conjecture whose themes are complementary: interpolation determinants settle three
independent bases; Gelfond--Schneider excludes algebraic irrational exponents; exact
certificates create a growing zero-free region; unique factorization bounds multiplicative
rank and now yields a common odd core; finite differences force local additive sparsity;
descent constrains perfect-power structure; and unconditional equivalences show that the
fractional-part axis measures the entire remaining distance to the conjecture.

The conjecture does not fall to a combination of these results, and the reason can now be
stated exactly.  A counterexample would not generate a messy exceptional set that density or
semigroup arguments might squeeze; it would generate the perfectly rigid free monoid
\[
\Sol=\Nzero+\Nzero\beta,\qquad\Cand=\{2^nM^k:n,k\in\Nzero\},
\]
whose block counts, fractional parts, and multiplicative relations are explicit and
internally consistent with every verified gap and density bound.  The remaining assertion is
that the second generator cannot exist --- equivalently
(Conjecture~\ref{conj:radical-target}), that no odd $u>1$ has $u^{\log_23}\in\Zthree$.

That reformulation is progress rather than pessimism: it removes artificial degrees of
freedom, supplies a canonical primitive object, upgrades density claims to exact theorems,
and gives a focused interface for future transcendence methods.  The near-term program is
listed in \cref{sec:program}; a complete proof still requires a new idea at the
four-exponentials boundary, but the target is now narrow and the surrounding formal
infrastructure is strong enough to test such an idea precisely.

\appendix

\section{Status matrix for major statements}
\label{app:status}

{\small
\begin{longtable}{@{}p{0.38\textwidth}p{0.21\textwidth}p{0.34\textwidth}@{}}
\toprule
Statement & Status & Location \\
\midrule
\endfirsthead
\toprule
Statement & Status & Location \\
\midrule
\endhead
$2^x,3^x,5^x\in\Z\Rightarrow x\in\Z$ & \statuslean & \lean{IntegerExponent} \\
Nonintegral solution irrational, transcendental & \statuslean &
\lean{TwoBaseIntegerExponent}, \lean{Transcendence} \\
$2^x<256\Rightarrow x\in\Z$ & \statuslean & \lean{FiniteCheck256} \\
$2^x<4096\Rightarrow x\in\Z$\ \ ($x\ge12$) & \statuslean & \lean{FiniteCheck4096} \\
Rank of candidate exponent vectors $\le2$ & \statuslean & \lean{MultiplicativeRank} \\
Common odd core $2^iw^j$ (existence) & \statuslean & \lean{OddCore} \\
Counting $\clog_2N\cdot\clog_3N$; linear dyadic blocks & \statuslean & \lean{OddCore} \\
Density dichotomy; gap equivalences; infinitude & \statuslean &
\lean{FractionalPartDensity} \\
$3$-smooth classification of further bases & \statuslean & \lean{ThreeSmooth} \\
Affine descent, $r+8k\le x$ & \statuslean & \lean{ArithmeticRigidity} \\
Three-term AP obstruction under $h^5\le n$ & \statuslean & \lean{ZeroDensity} \\
Shrinking-target Diophantine bounds & \statuslean & \lean{ShrinkingTarget} \\
Gelfond--Schneider theorem & \statuslean & \lean{GelfondSchneider} \\
$4$EC $\Rightarrow$ conjecture & \statuslean & \lean{FourExponentialsReduction} \\
Primitive (canonical) odd core; uniqueness & \statusnew & Theorem~\ref{thm:odd-core} \\
Exact solution monoid $\Nzero+\Nzero\beta$ & \statuslean &
Theorem~\ref{thm:generator}; \lean{PrimitiveGenerator} \\
Exact block counts; cumulative $T^2/(2\beta)$ & \statusnew & Theorem~\ref{thm:block-count} \\
Blockwise Weyl equidistribution & \statusnew & Theorem~\ref{thm:block-equidist} \\
Primitivity of least output pair & \statusnew &
Corollaries~\ref{cor:no-depth}, \ref{cor:no-common-power} and~\ref{cor:affine-primitive} \\
Sharp three-term AP obstruction ($h^{400}\le n^{83}$) & \statuslean &
\lean{SharpProgression} \\
Higher-order AP obstructions & \statusnew & Theorem~\ref{thm:higher-difference} \\
Localization reformulation & \statuslean &
Theorem~\ref{thm:localization}; \lean{Localization} \\
Radical-degree proposition & \statusnew & Proposition~\ref{prop:radical-degree} \\
No nonintegral candidate with $2^x<2^{20}$ & \statuscomp & Proposition~\ref{prop:finite20} \\
Full Alaoglu--Erd\H{o}s conjecture & open & reduced to Conjecture~\ref{conj:radical-target} \\
\bottomrule
\end{longtable}
}

\section{Formalization inventory}
\label{app:inventory}

All modules live under
\lean{Analysis/ExponentialIdentities/Lean/ExponentialIdentities/}, in namespace
\lean{LeanProofs}.  No \lean{sorry}, no extra axioms, no \lean{native_decide}; certificate
tables are replayed by the kernel.

\begin{center}
\small
\begin{tabular}{@{}ll@{}}
\toprule
Module & Contents \\
\midrule
\lean{IntegerExponent}(+3 submodules) & three-base theorem via interpolation determinants \\
\lean{TwoBaseIntegerExponent} & conjecture statement, candidates, gaps, $2^x<10$ \\
\lean{.../FiniteCheck}, \lean{FiniteCheck256} & zero-free regions $2^x<100$, $2^x<256$ \\
\lean{.../FiniteCheck4096} & zero-free region $2^x<4096$ (kernel \lean{decide} table) \\
\lean{.../ZeroDensity} & AP obstruction, Roth-number counting, $o(N)$ \\
\lean{.../NonintegerDensity} & nonintegral candidate refinement, density zero \\
\lean{.../MultiplicativeRank} & rank $\le2$, $\thetaq$ irrational, $(\clog_2N)^2$ \\
\lean{.../CandidateStructure} & cross-valuation identity, conditional box growth \\
\lean{.../OddCore} & common odd core, $\clog_2N\,\clog_3N$, dyadic blocks \\
\lean{.../ArithmeticRigidity} & descent, $r+8k\le x$ \\
\lean{.../ThreeSmooth} & $3$-smooth classification of further bases \\
\lean{.../ShrinkingTarget} & fract shrinking targets, Diophantine lower bounds \\
\lean{.../Transcendence} & GS applications: transcendence of solutions, of $\log_23$ \\
\lean{.../FractionalPartDensity} & density dichotomy, gap equivalences, infinitude \\
\lean{.../Localization} & conjecture $\iff$ no odd $u>1$ with $u^{\thetaq}\in\Zthree$ \\
\lean{.../SharpProgression} & sharp three-term AP obstruction, $\thetaq<317/200$ \\
\lean{.../PrimitiveGenerator} & conditional exact generator: $\Sol=\Nzero+\Nzero\beta$ \\
\lean{.../FourExponentialsReduction} & $4$EC $\Rightarrow$ conjecture \\
\lean{GelfondSchneider/...} & Gelfond--Schneider theorem (port of Karatarakis) \\
\bottomrule
\end{tabular}
\end{center}

\section{A symmetric formulation}
\label{app:symmetric}

Everything above has a base-swapped version.  Put $\eta=\log_32=1/\thetaq$.  Failure of the
conjecture is also equivalent to the existence of an integer $v>1$ with $3\nmid v$ and
\[
v^{\eta}\in\Z[1/2].
\]
For the primitive pair $(M,A)$, factoring $A$ into a power of $3$ times a $3$-coprime part
produces a symmetric primitive core and rational-power index; the two normalizations give
the same least exponent $\beta$, and their compatibility condition is exactly
$\log M\log3=\log A\log2$.  The symmetry is useful in formalization --- some divisibility
arguments are easier on one side --- and confirms that Corollary~\ref{cor:no-depth} is not an
artifact of privileging the base $2$.

\section{Interval computation: script and outputs}
\label{app:interval-output}

\begin{lstlisting}[style=pythonstyle,caption={Interval verification for one range.  Set START to 1 or $2^{19}$ and LIMIT to $2^{19}$ or $2^{20}$.},label={lst:interval}]
import math, time, hashlib
import mpmath as mp

START = 1
LIMIT = 1 << 19
DPS = 40

mp.iv.dps = DPS
ln2 = mp.iv.log(2)
ln3 = mp.iv.log(3)
theta_float = math.log(3.0) / math.log(2.0)

checked = 0
failures = []
min_lower = (float("inf"), None, None)
min_upper = (float("inf"), None, None)
manifest = hashlib.sha256()
start_time = time.time()

for m in range(START, LIMIT):
    if (m & (m - 1)) == 0:      # integral exponent: m is a power of two
        continue

    # Only proposes a neighboring integer.  The interval tests below certify it.
    n = math.floor(math.exp(theta_float * math.log(m)))

    lhs = mp.iv.log(m) * ln3
    lower_diff = lhs - mp.iv.log(n) * ln2
    upper_diff = mp.iv.log(n + 1) * ln2 - lhs

    lower_lb = float(lower_diff.a)
    upper_lb = float(upper_diff.a)
    if not (lower_lb > 0.0 and upper_lb > 0.0):
        failures.append((m, n, str(lower_diff), str(upper_diff)))
        break

    min_lower = min(min_lower, (lower_lb, m, n))
    min_upper = min(min_upper, (upper_lb, m, n))
    manifest.update(f"{m}:{n};".encode())
    checked += 1

print("checked_nonpowers=", checked)
print("failures=", len(failures))
print("min_lower_log_margin=", min_lower)
print("min_upper_log_margin=", min_upper)
print("sha256_m_n_manifest=", manifest.hexdigest())
print("elapsed_seconds=", time.time() - start_time)
assert not failures
\end{lstlisting}

The two 40-digit runs produced:

\begin{lstlisting}[style=pythonstyle,numbers=none,caption={Lower half output.}]
LIMIT=524288
DPS=40
checked_nonpowers=524268
failures=0
min_lower_log_margin=(1.4147089634746718e-15, 189427, 231498127)
min_upper_log_margin=(3.0841010256409534e-15, 306982, 497586055)
sha256_m_n_manifest=02b589461fa8c1721bc251575b98b1e3d58001f54cb7faa2cd6c59edb8c819a5
elapsed_seconds=18.505
\end{lstlisting}

\begin{lstlisting}[style=pythonstyle,numbers=none,caption={Upper half output.}]
LIMIT=1048576
DPS=40
checked_nonpowers=524287
failures=0
min_lower_log_margin=(9.490232037370244e-16, 958460, 3023921230)
min_upper_log_margin=(2.3019781003173366e-15, 556061, 1275861767)
sha256_m_n_manifest=fccccf5de8673379d37ff0fc840ab091d9ff7a2352ff64a93091f01c21bf4cb4
elapsed_seconds=18.904
\end{lstlisting}

Elapsed times are machine-dependent sanity checks; the counts, zero-failure results,
extremal tuples, and hashes are the reproducibility data.  The hashes cover the ASCII stream
\lean{m:n;} of proposed-and-certified pairs and make accidental changes to the candidate
sequence detectable across reruns.

\begin{thebibliography}{99}

\bibitem{AE1944}
L.~Alaoglu and P.~Erd\H{o}s,
\newblock \emph{On highly composite and similar numbers},
\newblock Transactions of the American Mathematical Society \textbf{56} (1944), 448--469.
\newblock DOI: \href{https://doi.org/10.1090/S0002-9947-1944-0011087-2}{10.1090/S0002-9947-1944-0011087-2}.

\bibitem{Ram1915}
S.~Ramanujan,
\newblock \emph{Highly composite numbers},
\newblock Proceedings of the London Mathematical Society (2) \textbf{14} (1915), 347--409;
annotated manuscript in Ramanujan Journal \textbf{1} (1997), 119--153.

\bibitem{Baker1975}
A.~Baker,
\newblock \emph{Transcendental Number Theory},
\newblock Cambridge University Press, 1975.

\bibitem{Hua1982}
L.-K.~Hua,
\newblock \emph{Introduction to Number Theory},
\newblock Springer, 1982, Chapter~17.9.

\bibitem{KuipersNiederreiter}
L.~Kuipers and H.~Niederreiter,
\newblock \emph{Uniform Distribution of Sequences},
\newblock Wiley, 1974; Dover reprint, 2006.

\bibitem{Roy1992}
D.~Roy,
\newblock \emph{Matrices whose coefficients are linear forms in logarithms},
\newblock Journal of Number Theory \textbf{41}(1) (1992), 22--47.
\newblock DOI: \href{https://doi.org/10.1016/0022-314X(92)90081-Y}{10.1016/0022-314X(92)90081-Y}.

\bibitem{Schneider1957}
Th.~Schneider,
\newblock \emph{Einf\"uhrung in die transzendenten Zahlen},
\newblock Springer, 1957.

\bibitem{Waldschmidt2000}
M.~Waldschmidt,
\newblock \emph{Diophantine Approximation on Linear Algebraic Groups},
\newblock Grundlehren der mathematischen Wissenschaften 326, Springer, 2000.
\newblock DOI: \href{https://doi.org/10.1007/978-3-662-11569-5}{10.1007/978-3-662-11569-5}.

\bibitem{WaldschmidtSix}
M.~Waldschmidt,
\newblock \emph{Six exponentials theorem --- irrationality},
\newblock expository note.

\bibitem{WaldschmidtVariations}
M.~Waldschmidt,
\newblock \emph{Variations on the Six Exponentials Theorem},
\newblock in \emph{Algebra and Number Theory}, pp.~338--355.
\newblock DOI: \href{https://doi.org/10.1007/978-93-86279-23-1_23}{10.1007/978-93-86279-23-1\_23}.

\bibitem{Waldschmidt2023}
M.~Waldschmidt,
\newblock \emph{The Four Exponentials Problem and Schanuel's Conjecture},
\newblock in \emph{Mathematics Going Forward}, Lecture Notes in Mathematics, 2023,
pp.~579--592.
\newblock DOI: \href{https://doi.org/10.1007/978-3-031-12244-6_39}{10.1007/978-3-031-12244-6\_39}.

\bibitem{TheoremDB2026}
TheoremDB contributors,
\newblock \emph{Four exponentials conjecture}, problem P4, reviewed July 31, 2026.
\newblock \url{https://theoremdb.org/statements/P4/}.

\bibitem{ProveItRepo}
\emph{ProveIt repository},
\newblock \url{https://github.com/VladimirReshetnikov/ProveIt}.

\bibitem{mpmath}
F.~Johansson et al.,
\newblock \emph{mpmath: a Python library for arbitrary-precision floating-point arithmetic},
\newblock version 1.3.0.
\newblock \url{https://mpmath.org/}.

\end{thebibliography}

\end{document}
